\documentclass[aspectratio=169,10pt]{beamer}
\usepackage{srlecture}
% Shared settings for the three QIPPG103 decks. Change the game URL here once.
\newcommand{\gameurl}{https://ebinrajiiit.github.io/ml_games/}
\newcommand{\gameurlshort}{ebinrajiiit.github.io/ml\_games}
\author{Dr.\ Sara Renjit}
\institute{Indian Institute of Information Technology Kottayam}
\newcommand{\coursename}{QIPPG103 Machine Learning: Principles and Practices}

\title{Reinforcement Learning: Formulating the Task}
\renewcommand{\lecturelabel}{QIPPG103 \textperiodcentered\ Lecture 24 \textperiodcentered\ Reinforcement Learning}

% campus grid: \campusgrid{path coordinates as (col,row) list}{colour}
\newcommand{\campusgrid}[2]{%
  \begin{tikzpicture}[x=0.46cm,y=-0.46cm]
    \foreach \c in {0,...,4}{\foreach \r in {0,...,4}{\draw[rulegrey,fill=white] (\c,\r) rectangle ++(1,1);}}
    \foreach \c in {1,...,4}{\fill[navy!22] (\c,2) rectangle ++(1,1); \node[font=\tiny,text=navy] at (\c+0.5,2.5) {$\sim$};}
    \fill[leaf!25] (2,0) rectangle ++(1,1); \node[font=\scriptsize\bfseries] at (2.5,0.5) {L};
    \node[font=\scriptsize\bfseries] at (2.5,4.5) {H};
    \draw[rulegrey!50!black] (0,0) rectangle (5,5);
    \draw[#2,line width=1.4pt,-{Stealth[length=2mm]},rounded corners=2pt] #1;
  \end{tikzpicture}}

\begin{document}

\begin{frame}[plain]
  \begin{tikzpicture}[remember picture,overlay]
    \node[anchor=north east] at ([xshift=-1.2em,yshift=-1.0em]current page.north east) {\includegraphics[height=1.25cm]{iiitk.png}};
  \end{tikzpicture}
  \vskip3.2em
  {\small\scshape\color{inksoft}\coursename\par}
  \vskip1.2em
  {\color{navy}\fontsize{28}{32}\selectfont Reinforcement Learning\par}
  \vskip0.3em
  {\color{navy}\Large Applications, task formulation, and defining an environment\par}
  \vskip0.5em
  {\large\itshape Learning from consequences instead of answers\par}
  \vskip0.6em{\color{rulegrey}\rule{0.45\textwidth}{0.4pt}}\par\vskip0.8em
  {\normalsize Dr.\ Sara Renjit\par}
  {\small\color{inksoft}Indian Institute of Information Technology Kottayam\par}
  \vskip0.4em
  {\small\color{inksoft}Lecture 24 of 27 \textperiodcentered\ offline session \textperiodcentered\ 90 minutes\par}
\end{frame}

\begin{frame}{How did you learn to ride a bicycle?}
  \vskip0.6em
  \begin{columns}[T]
    \begin{column}{0.5\textwidth}
      Nobody gave you a dataset of\\[0.3em]
      \quad (angle of lean, speed) $\;\rightarrow\;$ correct handle turn.
      \vskip0.8em
      You tried something, wobbled, fell, tried again. Staying up felt good; falling hurt. After a few days your body had a rule for what to do in every situation.
    \end{column}
    \begin{column}{0.46\textwidth}
      \uncover<2->{
      That is the whole idea of reinforcement learning:
      \begin{itemize}
        \item no correct answer is given for each situation;
        \item you act, and the world answers with a \emph{consequence};
        \item consequences can arrive late (you fall three seconds after the bad turn);
        \item what you do changes what you see next.
      \end{itemize}}
    \end{column}
  \end{columns}
\end{frame}

\begin{frame}{Where reinforcement learning sits in this course}
  \small\renewcommand{\arraystretch}{1.2}
  \begin{tabular}{@{}>{\raggedright}p{2.7cm}>{\raggedright}p{3.6cm}>{\raggedright}p{3.6cm}>{\raggedright\arraybackslash}p{3.7cm}@{}}
    \toprule
    & \textbf{Supervised} (Lect.\ 3--12) & \textbf{Unsupervised} (Lect.\ 13--17) & \textbf{Reinforcement} (Lect.\ 24--26) \\ \midrule
    Training signal & the correct label for every input & none & a number, the reward, after each action \\
    \pause
    Feedback & immediate and exact & -- & delayed, and only says ``how good'', never ``what was right'' \\
    \pause
    Data & fixed dataset, examples independent & fixed dataset & generated by the learner's own actions; one step depends on the last \\
    \pause
    Goal & predict labels & find structure & choose actions that maximise total future reward \\ \bottomrule
  \end{tabular}
  \vskip0.8em
  \small The last two rows are what make RL hard: the learner must explore to find good data, and must work out which of many earlier actions earned a reward.
\end{frame}

\begin{frame}{Where it has worked}
  \small\renewcommand{\arraystretch}{1.15}
  \begin{tabular}{@{}>{\raggedright}p{3.1cm}>{\raggedright}p{5.5cm}>{\raggedright\arraybackslash}p{5.0cm}@{}}
    \toprule
    \textbf{Application} & \textbf{What the agent controls} & \textbf{Reward} \\ \midrule
    Board games\newline {\scriptsize AlphaGo, 2016} & which stone to place & $+1$ win, $-1$ loss, at the very end \\
    \pause
    Data-centre cooling\newline {\scriptsize Google DeepMind, 2016} & fan, chiller and pump settings & minus energy used, within safe temperatures \\
    \pause
    Language models\newline {\scriptsize RLHF, 2022 onward} & the next token of a reply & a score from a model trained on human preferences \\
    \pause
    Traffic signals & which phase is green and for how long & minus vehicles waiting \\
    \pause
    Robots & joint torques & progress towards the goal, minus falls \\
    \pause
    Recommendations & which item to show next & clicks, watch time, return visits \\ \bottomrule
  \end{tabular}
  \vskip0.6em
  \small In every row someone first had to decide what the agent sees, what it may do and what counts as good. That decision is today's topic.
\end{frame}

\sectionslide{Part 1}{The vocabulary}

\begin{frame}{The agent--environment loop}
  \begin{columns}[c]
    \begin{column}{0.5\textwidth}
      \centering
      \begin{tikzpicture}[box/.style={draw=navy,line width=1pt,rounded corners=3pt,minimum width=3.4cm,minimum height=1.0cm,font=\normalsize}]
        \node[box] (ag) at (0,2.6) {Agent};
        \node[box] (en) at (0,0) {Environment};
        \draw[arr,visible on=<2->] (ag.east) -- (2.6,2.6) -- node[right,font=\small]{action $a_t$} (2.6,0) -- (en.east);
        \draw[arr,visible on=<3->] ([yshift=-0.25cm]en.west) -- (-2.9,-0.25) -- node[left,font=\small]{state $s_{t+1}$} (-2.9,2.85) -- ([yshift=0.25cm]ag.west);
        \draw[rec,visible on=<4->] ([yshift=0.25cm]en.west) -- (-2.2,0.25) -- node[right,font=\small,text=brick]{reward $r_{t+1}$} (-2.2,2.35) -- ([yshift=-0.25cm]ag.west);
      \end{tikzpicture}
    \end{column}
    \begin{column}{0.48\textwidth}
      At each time step $t = 0,1,2,\ldots$
      \begin{enumerate}[<+(1)->]
        \item the agent observes state $s_t$;
        \item it picks action $a_t$;
        \item the environment moves to $s_{t+1}$ and hands back a reward $r_{t+1}$.
      \end{enumerate}
      \vskip0.6em
      This produces a trajectory
      \[ s_0,\,a_0,\,r_1,\,s_1,\,a_1,\,r_2,\,s_2,\,\ldots \]
      Note the parallel with Lecture 20: a sequence, processed one step at a time.
    \end{column}
  \end{columns}
\end{frame}

\begin{frame}{Seven words to know}
  \small\renewcommand{\arraystretch}{1.2}
  \begin{tabular}{@{}>{\raggedright}p{2.6cm}>{\raggedright}p{6.0cm}>{\raggedright\arraybackslash}p{5.0cm}@{}}
    \toprule
    \textbf{Term} & \textbf{Meaning} & \textbf{Bicycle} \\ \midrule
    State $s$ & what the agent knows about the situation now & speed, lean angle, road ahead \\
    \pause
    Action $a$ & a choice it can make & turn the handle, pedal, brake \\
    \pause
    Reward $r$ & a number saying how good the last step was & $0$ while upright, $-10$ for a fall \\
    \pause
    Policy $\pi(a\mid s)$ & its rule for choosing actions & the habit in your arms \\
    \pause
    Episode & one run from start to an end state & one ride, home to college \\
    \pause
    Return $G_t$ & total (discounted) reward from $t$ onward & how the whole ride went \\
    \pause
    Discount $\gamma$ & how much a later reward counts, $0\le\gamma\le1$ & do you care about tomorrow's sore knee? \\ \bottomrule
  \end{tabular}
  \vskip0.6em
  \small The agent's goal: find the policy $\pi$ that makes the expected return as large as possible.
\end{frame}

\begin{frame}{The return, and why we discount}
  \[
    G_t \;=\; r_{t+1} + \gamma\, r_{t+2} + \gamma^2 r_{t+3} + \cdots \;=\; \sum_{k=0}^{\infty} \gamma^k\, r_{t+k+1}
  \]
  \begin{columns}[T]
    \begin{column}{0.5\textwidth}
      \textbf{Worked example}, $\gamma = 0.9$\\[0.3em]
      Rewards along an episode: $-1,\;-1,\;-1,\;+10$
      \vskip0.3em
      \renewcommand{\arraystretch}{1.2}
      \begin{tabular}{@{}l@{}}
        $G_0 = -1 + 0.9(-1) + 0.81(-1) + 0.729(10)$ \\
        \uncover<2->{$\phantom{G_0} = -1 - 0.9 - 0.81 + 7.29$} \\
        \uncover<3->{$\phantom{G_0} = \alert{4.58}$}
      \end{tabular}
    \end{column}
    \begin{column}{0.47\textwidth}
      \textbf{Why $\gamma<1$?}
      \begin{itemize}
        \item In a task that never ends, an undiscounted sum can be infinite.
        \item A reward now is more certain than one later.
        \item It makes the agent prefer the \emph{shorter} of two routes to the same prize.
      \end{itemize}
      \vskip0.3em
      $\gamma=0$: only the next reward matters (short-sighted).\\
      $\gamma\to1$: all future rewards count almost equally.
    \end{column}
  \end{columns}
\end{frame}

\begin{frame}{The Markov property: is your state enough?}
  \begin{columns}[T]
    \begin{column}{0.52\textwidth}
      A state is \textbf{Markov} if the future depends only on it, not on how you got there:
      \[ P(s_{t+1}, r_{t+1}\mid s_t, a_t) = P(s_{t+1}, r_{t+1}\mid s_0,a_0,\ldots,s_t,a_t) \]
      With Markov states and actions, the task is a \textbf{Markov decision process} (MDP):
      \[ (\mathcal S,\ \mathcal A,\ P,\ R,\ \gamma) \]
    \end{column}
    \begin{column}{0.45\textwidth}
      \textbf{Quick check}
      \begin{itemize}
        \item Chess: the board position (plus whose turn, castling rights). Markov.
        \item Pong from \emph{one} screenshot: you see the ball, not where it is going. Not Markov.\\ Fix: stack the last 4 frames (this is what the Atari DQN did).
        \item A patient's vitals right now, for a dosing agent: probably not; add recent history.
      \end{itemize}
    \end{column}
  \end{columns}
  \vskip0.6em
  \margin{Link to Lectures 20--21: when the raw observation is not Markov, an RNN or LSTM can summarise the history into a state.}
\end{frame}

\sectionslide{Part 2}{Formulating a task: five questions}

\begin{frame}{A recipe for turning a problem into an RL task}
  \begin{enumerate}
    \item \textbf{Who is the agent, and what is the environment?} Draw the boundary. Anything the agent cannot change directly is environment.
    \item \textbf{State space $\mathcal S$.} What must the agent see to choose well? Is it Markov? Discrete or continuous? How big?
    \item \textbf{Action space $\mathcal A$.} Discrete (left/right/stay) or continuous (steering angle in $[-30^\circ,30^\circ]$)? Are some actions illegal in some states?
    \item \textbf{Reward $R$.} What number, when? Reward the \emph{outcome you want}, not the method you guess will get there.
    \item \textbf{Episodes.} Is there a natural end (episodic: a game, a delivery) or does it run for ever (continuing: a thermostat)? What is $\gamma$? Is there a time limit?
  \end{enumerate}
  \vskip0.6em
  \small Sutton and Barto, Chapter 3, is the reference for this lecture.
\end{frame}

\begin{frame}{Four tasks, formulated}
  \small
  \renewcommand{\arraystretch}{1.2}
  \begin{tabular}{@{}>{\raggedright}p{1.8cm}>{\raggedright}p{3.6cm}>{\raggedright}p{2.7cm}>{\raggedright}p{3.0cm}>{\raggedright\arraybackslash}p{2.0cm}@{}}
    \toprule
    & \textbf{State} & \textbf{Actions} & \textbf{Reward} & \textbf{Episodes} \\ \midrule
    Chess & board, side to move, castling and en passant rights & legal moves (discrete, varies by state) & $+1$ / $0$ / $-1$ at the end only & episodic, $\gamma=1$ \\
    \pause
    Room thermostat & room temp., outside temp., time of day, occupancy & heater power, 0--100\% (continuous) & $-$energy $-\lambda\,|$temp $-$ 24$^\circ$C$|$ & continuing, $\gamma=0.99$ \\
    \pause
    Junction signal & queue length on each arm, current phase, time in phase & keep phase, or switch (discrete) & $-$(vehicles waiting) each step & continuing \\
    \pause
    Robot vacuum & position, map of dirty tiles, battery & move N/E/S/W, dock & $+1$ per tile cleaned, $-100$ if battery dies & episodic \\ \bottomrule
  \end{tabular}
  \vskip0.6em
  Notice: the thermostat's reward \emph{trades off} two goals with a weight $\lambda$. Choosing $\lambda$ is a design decision, not something the algorithm learns.
\end{frame}

\begin{frame}{When the reward says the wrong thing}
  \begin{columns}[T]
    \begin{column}{0.52\textwidth}
      An agent maximises exactly the number you give it, and nothing else.
      \vskip0.5em
      \textbf{A real case.} OpenAI trained an agent on the boat-racing game \emph{CoastRunners}, rewarding the in-game score. The agent found a lagoon where three targets kept respawning, and circled there for ever, on fire, crashing into boats, never finishing the race. It scored higher than human players (Clark and Amodei, 2016).
      \vskip0.5em
      This is called \textbf{reward hacking} or \textbf{specification gaming}.
    \end{column}
    \begin{column}{0.45\textwidth}
      \textbf{Common mistakes}
      \begin{itemize}
        \item Rewarding activity (``+1 per step working'') instead of results.
        \item Rewarding a sub-goal so richly that it beats the goal.
        \item Forgetting a cost (energy, wear, safety) the designer cared about but never wrote down.
        \item Equal rewards for slow and fast success, with $\gamma=1$.
      \end{itemize}
      \vskip0.4em
      You will make all four on purpose in today's game.
    \end{column}
  \end{columns}
\end{frame}

\begin{frame}{Worksheet, Part B: formulate one yourself}
  \yourturn{10 min}
  \begin{block}{In groups of three, choose one and fill in the four-row table}
    \begin{enumerate}[<1->]
      \item[A.] The Kottayam KSRTC bus stand: an agent decides which bay each arriving bus is sent to.
      \item[B.] An irrigation pump for a paddy field in Kuttanad: when to pump and for how long.
      \item[C.] A tutoring app that picks the next practice question for a student.
    \end{enumerate}
  \end{block}
  \vskip0.6em
  Then write one reward that a careless designer might choose, and say how the agent could exploit it. Two groups will present in 60 seconds each.
\end{frame}

\sectionslide{Part 3 \textperiodcentered\ case study}{Campus Courier: defining an environment}

\begin{frame}{The task}
  \begin{columns}[c]
    \begin{column}{0.36\textwidth}
      \centering
      \begin{tikzpicture}[x=0.85cm,y=-0.85cm]
        \foreach \c in {0,...,4}{\foreach \r in {0,...,4}{\draw[rulegrey,fill=white] (\c,\r) rectangle ++(1,1);}}
        \foreach \c in {1,...,4}{\fill[navy!22] (\c,2) rectangle ++(1,1); \node[text=navy] at (\c+0.5,2.5) {$\sim$};}
        \fill[leaf!25] (2,0) rectangle ++(1,1); \node[font=\bfseries] at (2.5,0.5) {L};
        \node[font=\bfseries] at (2.5,4.5) {H};
        \draw[rulegrey!50!black] (0,0) rectangle (5,5);
        \foreach \c in {0,...,4}{\node[font=\tiny,text=inksoft] at (\c+0.5,-0.25) {\c};}
        \foreach \r in {0,...,4}{\node[font=\tiny,text=inksoft] at (-0.25,\r+0.5) {\r};}
      \end{tikzpicture}\\[0.3em]
      {\scriptsize H hostel (start) \quad L library (goal)\\ $\sim$ rain puddle, dry gap at column 0}
    \end{column}
    \begin{column}{0.6\textwidth}
      A delivery robot must carry books from the hostel to the library after the rain.
      \vskip0.5em
      \renewcommand{\arraystretch}{1.25}
      \begin{tabular}{@{}>{\raggedright}p{1.6cm}>{\raggedright\arraybackslash}p{6.4cm}@{}}
        Agent & the robot \\
        State & its cell: 25 states, numbered $5\cdot\text{row}+\text{col}$ \\
        Actions & up, right, down, left (4, discrete); hitting the edge means staying put \\
        Reward & step $-1$, entering a puddle $-10$, reaching L $+10$ \\
        Episodes & ends at L (terminated) or after 30 moves (truncated); $\gamma=0.95$ \\
      \end{tabular}
      \vskip0.5em
      \small Straight through the puddle: 4 moves. The dry way round: 8 moves.
    \end{column}
  \end{columns}
\end{frame}

\begin{frame}[fragile]{Defining it in code: the Gymnasium interface}
\begin{lstlisting}
class CampusCourier(gym.Env):
    def __init__(self, step_reward=-1.0, puddle_reward=-10.0, goal_reward=10.0, max_steps=30):
        self.observation_space = spaces.Discrete(25)       # S: 25 cells
        self.action_space      = spaces.Discrete(4)        # A: up, right, down, left
        ...
    def reset(self, seed=None, options=None):
        super().reset(seed=seed)
        self.pos, self.t = (4, 2), 0                       # start at the hostel
        return self._obs(), {}

    def step(self, action):
        dr, dc = self.MOVES[action]
        r, c = self.pos[0] + dr, self.pos[1] + dc
        if 0 <= r < 5 and 0 <= c < 5: self.pos = (r, c)    # walls: stay put
        self.t += 1
        terminated = self.pos == self.GOAL                 # a real end
        truncated  = self.t >= self.max_steps              # an artificial time limit
        reward = (self.goal_reward if terminated else
                  self.puddle_reward if self.pos in self.PUDDLE else self.step_reward)
        return self._obs(), reward, terminated, truncated, {}
\end{lstlisting}
  \small Every RL library expects exactly these two methods. \margin{Full file, with a test: \texttt{case\_study/campus\_courier\_env.py}. It passes Gymnasium's \texttt{check\_env}.}
\end{frame}

\begin{frame}{First test: a robot that moves at random}
  \begin{columns}[T]
    \begin{column}{0.5\textwidth}
      Before any learning, always run a random policy. It checks the environment works and gives a baseline.
      \vskip0.6em
      \textbf{1{,}000 random episodes} (from the script):
      \begin{itemize}
        \item reached the library in \alert{27.5\%} of them;
        \item average discounted return \alert{$-32.28$}.
      \end{itemize}
    \end{column}
    \begin{column}{0.47\textwidth}
      \textbf{Two hand-written routes}, $\gamma=0.95$
      \vskip0.3em
      \renewcommand{\arraystretch}{1.25}
      \begin{tabular}{@{}lrr@{}}
        \toprule
        Route & moves & return \\ \midrule
        through the puddle & 4 & $-2.83$ \\
        \pause
    dry way round & 8 & $0.95$ \\ \bottomrule
      \end{tabular}
      \vskip0.6em
      With this reward the longer dry route is worth more. Change the puddle penalty and that can flip. The reward \emph{is} the specification.
    \end{column}
  \end{columns}
\end{frame}

\begin{frame}{Same robot, three reward designs}
  \small
  Best policy for each design, computed exactly by value iteration ($\gamma=0.95$); learning it by trial and error is Q-learning, Lecture 25.
  \vskip0.2em
  \begin{columns}[T]
    \begin{column}{0.32\textwidth}
      \centering
      \campusgrid{(2.5,4.5) -- (2.5,3.5) -- (4.5,3.5) -- (4.85,3.5)}{brick}\\[0.3em]
      \textbf{step $+1$}, puddle 0, goal $+10$\\[0.2em]
      \raggedright Pushes against a wall for ever: $+1$ a step is worth $1/(1-0.95)=20$, beating $+10$.
    \end{column}
    \begin{column}{0.32\textwidth}\uncover<2->{
      \centering
      \campusgrid{(2.5,4.5) -- (2.5,0.5)}{ochre}\\[0.3em]
      \textbf{step $-1$}, puddle $-1$, goal $+10$\\[0.2em]
      \raggedright Goes straight through the puddle. Nothing in the reward says wet books are bad.
    }\end{column}
    \begin{column}{0.32\textwidth}\uncover<3->{
      \centering
      \campusgrid{(2.5,4.5) -- (2.5,3.5) -- (0.5,3.5) -- (0.5,0.5) -- (2.5,0.5)}{leaf}\\[0.3em]
      \textbf{step $-1$}, puddle $-10$, goal $+10$\\[0.2em]
      \raggedright Takes the dry route, 8 moves. The puddle penalty must beat the cost of 4 extra steps.
    }\end{column}
  \end{columns}
\end{frame}

\begin{frame}{Game: Reward Designer}
  \yourturn{12 min}
  \begin{columns}[c]
    \begin{column}{0.36\textwidth}
      \centering
      \qrcode[height=4.0cm]{\gameurl reward-designer.html}\\[0.4em]
      {\scriptsize\ttfamily \gameurlshort}
    \end{column}
    \begin{column}{0.6\textwidth}
      \textbf{Round 1: Formulate it.} Four situations. Pick the right state, action, reward or episode type.
      \vskip0.5em
      \textbf{Round 2: Fix the reward.} Three broken designs for the Campus Courier. Move the sliders for step, puddle and goal rewards and $\gamma$, press \emph{Train}, and watch where the robot goes.
      \vskip0.5em
      Record your final slider values for each challenge in worksheet Part C. We compare across the room afterwards.
    \end{column}
  \end{columns}
\end{frame}

\begin{frame}{What the game was teaching}
  \begin{itemize}
    \item \textbf{Challenge 1.} A positive reward per step makes staying alive the goal. Any step reward $\le 0$ fixed it. So did a goal reward well above 20, the value of loitering for ever at $+1$ a step.
    \item \textbf{Challenge 2.} The robot gets wet unless the puddle penalty outweighs the extra moves. With step $-1$ and $\gamma=0.95$ the switch is at a puddle reward of about $-6.02$: solve $6.67+0.95\,r_p = 0.95$.
    \item \textbf{Challenge 3.} With no step cost and $\gamma=1$, every route to the library is worth the same $+10$, so the robot has no reason to arrive. Any $\gamma<1$ makes sooner better than later.
  \end{itemize}
  \vskip0.6em
  \begin{block}{The point}
    The algorithm was the same in every case. Only the formulation changed, and it decided everything.
  \end{block}
\end{frame}

\begin{frame}{Exit ticket (worksheet Part D)}
  \yourturn{5 min}
  \begin{enumerate}
    \item Rewards $0, 0, 0, 1$ with $\gamma=0.5$. What is $G_0$?\quad (a) 1 \quad (b) 0.5 \quad (c) 0.125 \quad (d) 0.0625
    \item A self-driving car sets its steering angle anywhere in $[-30^\circ, 30^\circ]$. Its action space is\\ (a) discrete \quad (b) continuous \quad (c) Markov \quad (d) episodic
    \item An agent sees a single frame of a ball game and cannot tell which way the ball is moving. The problem is that its state is\\ (a) not Markov \quad (b) too large \quad (c) continuous \quad (d) not discounted
    \item Which reward for a warehouse robot invites reward hacking?\\ (a) $+1$ per parcel delivered \quad (b) $+1$ per metre travelled \quad (c) $-1$ per minute late \quad (d) $-5$ per collision
  \end{enumerate}
  \vskip0.4em
  \uncover<+->{\small Answers: 1 (c) $0.5^3=0.125$ \quad 2 (b) \quad 3 (a) \quad 4 (b): it can drive in circles.}
\end{frame}

\begin{frame}{To take away}
  \begin{enumerate}
    \item RL learns from rewards for actions, not from labelled answers; the feedback is delayed and the data depends on the learner's own choices.
    \item Formulating a task means fixing the agent, the state space, the action space, the reward, and the episode structure with $\gamma$.
    \item The return $G_t=\sum_k \gamma^k r_{t+k+1}$ is what the agent maximises; $\gamma$ sets how far ahead it looks.
    \item Check the state is Markov; if not, add history.
    \item The agent optimises exactly the reward you wrote. Test it with a random policy and with hand-written routes before training anything.
  \end{enumerate}
  \vskip0.8em
  \textbf{Next (Lecture 25, Dr Manu Madhavan):} function approximation and Q-learning, the algorithm that learns the Campus Courier policy without being given the map.
\end{frame}

\begin{frame}{References}
  \small
  \begin{itemize}[<1->]
    \item R.\ S.\ Sutton, A.\ G.\ Barto. \emph{Reinforcement Learning: An Introduction}, 2nd ed. MIT Press, 2018. Chapters 1 and 3.
    \item D.\ Silver et al. Mastering the game of Go with deep neural networks and tree search. \emph{Nature} 529:484--489, 2016.
    \item V.\ Mnih et al. Human-level control through deep reinforcement learning. \emph{Nature} 518:529--533, 2015.
    \item J.\ Clark, D.\ Amodei. Faulty reward functions in the wild. OpenAI blog, 2016.
    \item R.\ Evans, J.\ Gao. DeepMind AI reduces Google data centre cooling bill by 40\%. DeepMind blog, 2016.
    \item L.\ Ouyang et al. Training language models to follow instructions with human feedback. \emph{NeurIPS} 2022.
    \item Farama Foundation. Gymnasium documentation, \texttt{gymnasium.farama.org}.
  \end{itemize}
\end{frame}

\end{document}
